\documentclass{article}
\usepackage[english]{babel}
\usepackage[letterpaper,top=2cm,bottom=2cm,left=2.5cm,right=2.5cm,marginparwidth=1.75cm]{geometry}
\usepackage{amsmath}
\usepackage{graphicx}
\usepackage[colorlinks=true, allcolors=blue]{hyperref}
\usepackage[numbers]{natbib}
\usepackage[section]{placeins}

% Long \texttt object names (kNue2024FDAllSamples_Prod62024EffMatched etc.) cannot
% hyphenate, so give TeX room to stretch a line rather than run into the margin.
\setlength{\emergencystretch}{3em}

\setlength{\textfloatsep}{8pt plus 2pt minus 2pt}
\setlength{\floatsep}{8pt plus 2pt minus 2pt}
\setlength{\intextsep}{8pt plus 2pt minus 2pt}

\graphicspath{{images-ssahoo/}}

% ---------------------------------------------------------------------------
% Figure helpers. The plots are a selection from the talk (NOVA-doc-70698);
% that reference has the full set, covering every sample and both estimators.
% ---------------------------------------------------------------------------
% \twofig{a}{b}{caption}{label}          -- one row of two
\newcommand{\twofig}[4]{%
\begin{figure}[htbp]\centering
  \includegraphics[width=0.46\textwidth]{#1}\hfill
  \includegraphics[width=0.46\textwidth]{#2}
  \caption{#3}\label{#4}
\end{figure}}
% \threefig{a}{b}{c}{caption}{label}     -- one row of three
\newcommand{\threefig}[5]{%
\begin{figure}[htbp]\centering
  \includegraphics[width=0.315\textwidth]{#1}\hfill
  \includegraphics[width=0.315\textwidth]{#2}\hfill
  \includegraphics[width=0.315\textwidth]{#3}
  \caption{#4}\label{#5}
\end{figure}}
% \fourfig{tl}{tr}{bl}{br}{caption}{label} -- two rows of two
\newcommand{\fourfig}[6]{%
\begin{figure}[htbp]\centering
  \includegraphics[width=0.46\textwidth]{#1}\hfill
  \includegraphics[width=0.46\textwidth]{#2}\\[2pt]
  \includegraphics[width=0.46\textwidth]{#3}\hfill
  \includegraphics[width=0.46\textwidth]{#4}
  \caption{#5}\label{#6}
\end{figure}}

\title{NOvA Production 6 Validation Task Force Contribution Summary}
\author{Subhadarsi Sahoo}

\begin{document}
\maketitle

\section{Task Description}

My contribution to the Production~6 Validation Task Force was to validate the
\textbf{electron-neutrino energy estimators}: the \textbf{Spline}
(\texttt{kNueEnergy2024\_Prod6}), a quadratic fit in the EM and hadronic
calorimetric energies with separate FHC and RHC fits, and the \textbf{regression
CNN} (\texttt{regCVNNueE}, reading \texttt{energy.nue.regcvnEvtE}), an end-to-end
network trained on Prod6. Both are filled from the same selected events with the
same weight, so the comparison is unbiased by construction.

Eight Prod6 samples were processed: six \textbf{nominal} MC samples --- ND
FHC/RHC, and FD FHC/RHC non-swap and flux-swap --- together with ND FHC/RHC
\emph{data}. The MC is nominal throughout, with no systematically shifted
samples. The selection is the official Ana2024
Prod6 chain --- \texttt{kNue2024ND\_\allowbreak Prod62024EffMatched} at the ND and
\texttt{kNue2024FDAllSamples\_\allowbreak Prod62024EffMatched} (core $\cup$ peripheral) at the
FD --- with \texttt{kPPFXFluxCVWgt} on MC and \texttt{kUnitWeight} on data. No
cross-section CV tune is applied: none exists for Prod6
(\texttt{kXSecCVWgt2024} throws \texttt{Unsupported\allowbreak GeneratorException} on this
GENIE), so PPFX-only is the agreed recipe.

The benchmarks are energy spectra, resolution, bias and width
against true energy, migration matrices, the true-interaction-type composition
(including per PID bin), and ND data/MC. The work was reported in
Refs.~\cite{docdb-70195, docdb-69940, docdb-70459, docdb-70698}. The figures
below are a representative selection; the full set, covering every sample and
both estimators, is in Ref.~\cite{docdb-70698}.

\section{Findings}

\paragraph{Energy spectra, bias and resolution.}
Both estimators reconstruct the $\nu_e$ energy well in all six MC samples, and
their spectra and resolutions are reasonable throughout. The CNN has the narrower
core resolution in every sample. It is also the less biased of the two in five of
the six; the exception is FHC FD flux-swap, where both sit within a fraction of a
percent of zero and the Spline is nominally the closer --- so the CNN cannot
simply be described as less biased everywhere. The flux-swap samples resolve
better than the non-swap ones, as expected for a pure signal $\nu_e$ sample
compared with one containing NC $\pi^0$ background. Both estimators are near-unbiased at low energy and
increasingly under-reconstruct towards the top of the selection window; this is
an acceptance effect, since the window is applied in reconstructed energy, rather
than a failure of either estimator. Bias and width are quoted from an iterative
$\pm2\sigma$ core fit rather than a full-range fit, which is dragged by the
low-side tail and describes neither the core nor the tail; earlier talks quoted
the full-range RMS, which is the same distribution under a different summary
statistic and should not be compared with the core width.

\paragraph{True-interaction-type composition.}
The composition of the selected sample was determined for every sample, both
overall and split by PID bin. The ND and FD non-swap selections are roughly half
$\nu_e$~CC in FHC and closer to two thirds in RHC, with the remainder dominated
by NC $\pi^0$. RHC is the purer at both detectors because the antineutrino beam
carries far less $\nu_\mu$ to be mistaken for signal. The FD flux-swap samples
are almost entirely $\nu_e$~CC by construction. At the FD, purity rises with PID score in every
sample, with the peripheral sample falling between the two core bins. The truth
categories were defined to be exhaustive and mutually exclusive so that they sum
to unity by construction, which provides a built-in check on every composition
plot; the official \texttt{TruthCuts} were deliberately not used because they do
not form a partition.

\paragraph{ND data/MC.}
Data lies above MC in both beams, by roughly 20\% in FHC and 33\% in RHC, using
the same selection on data and MC with the MC POT-scaled to the data exposure.
The excess is beam-dependent, which disfavours a detector effect since that
would be common to both beams. Further study will depend on the cross-section
tuning once it exists for Prod6.

\paragraph{Issues found that affect other Prod6 users.}
\begin{itemize}
  \setlength{\itemsep}{1pt}
  \item \textbf{Most of the official \texttt{SharedPlotLists} $\nu_e$ variables do
  not work on Prod6.} Of the 63 registered, \textbf{49 place $100\%$ of events in
  the underflow bin}, in every sample at both detectors --- they are 2020-era
  \texttt{Var}s reading pre-Prod6 branches, and \texttt{sel.cvn2d} does not exist
  in these CAFs at all. They fail silently, so the plots are produced and look
  empty rather than raising an error. The remaining 14 read branches that do
  exist in Prod6 and are fine.
  \item \textbf{The stock \texttt{kLongestProng} returns 0 for every Prod6
  event}, because it reads the pre-Prod6 \texttt{vtx.elastic} path.
  \texttt{kLongestProng\_Prod6} must be used instead. A \texttt{Var} reading a
  missing branch does not error --- it returns 0 and the cut silently rejects
  everything, which is the common root cause of this item and the one above.
  \item \textbf{\texttt{regcvnEvtE\_opphorn} is bit-identical to
  \texttt{regcvnEvtE}} in Prod6 --- no bin of the two distributions differs. The
  branch duplicates the nominal rather than carrying an independently trained
  network, so the horn-training robustness check for the $\nu_e$ CNN cannot be
  performed with Prod6 CAFs; a naive comparison would report perfect robustness.
  Only the $\nu_e$ branch was checked.
  \item \textbf{There is no Prod6 $\nu_e$ analysis axis and no
  \texttt{Prod6Loaders.h}.} All four axis and analysis-bin symbols in
  \texttt{NueCuts2024.h} are pre-Prod6, and the pre-Prod6 axis cannot simply be
  reused on Prod6 files: it applies the original PID edges rather than the
  Prod6-retuned \texttt{\_Prod62024EffMatched} values, so events would be binned
  by a different selection than the one applied, silently. A Prod6 analysis bin
  with a matching axis, plus Prod6 loaders, would unblock Prod6 $\nu_e$
  predictions generally.
\end{itemize}

\paragraph{Future work.}
The Spline still carries the Ana2024 (Prod5.1) coefficients, so a good part of
its residual bias on Prod6 is simply those coefficients applied to a different
production. Retuning the polynomial on Prod6 is the immediate next task; it
should reduce that bias, and it will also make the comparison with the CNN a fair
one, since the CNN is already trained on Prod6. A Prod6 FD prediction using
\texttt{PredictionNoExtrap}, built from official CAFAna classes because the
official prediction chain cannot be pointed at Prod6, has been written and
validated but not yet run. Systematic uncertainties have not yet been studied
for these estimators.

\clearpage
\section{Near Detector}

\fourfig{energy_fhc_nd.png}{energy_rhc_nd.png}{res_fhc_nd.png}{res_rhc_nd.png}
{Reconstructed $\nu_e$ energy (top) and fractional resolution
$(E_{\rm reco}-E_{\rm true})/E_{\rm true}$ (bottom) at the ND, FHC left and RHC
right, for truth, Spline and CNN. RHC looks spikier only because it has a fifth
of the statistics; $\chi^2/$ndf $\approx1$ against a smooth reference.}
{fig:nd-energy-res}

\twofig{mean_fhc_nd.png}{std_fhc_nd.png}
{Bias (left) and width (right) of the fractional residual against true energy,
FHC ND. Both estimators are near-unbiased below $\sim$1.5~GeV and increasingly
under-reconstruct above $\sim$4~GeV, where the selection window truncates the
sample.}{fig:nd-bias-width}

\twofig{ts_sp_fhc_nd.png}{ts_sp_rhc_nd.png}
{ND composition by true interaction type, FHC left and RHC right, on the Spline
energy axis. The CNN gives an identical composition, as it must --- the same
events are selected, only the energy axis differs --- so only one estimator is
shown.}{fig:nd-composition}

\clearpage
\section{Near Detector data/MC}

\fourfig{dmc_spl_fhc.png}{dmc_spl_rhc.png}{dmc_cnn_fhc.png}{dmc_cnn_rhc.png}
{Reconstructed $\nu_e$ energy, data over MC, POT-scaled: Spline top, CNN bottom;
FHC left, RHC right.}{fig:nd-datamc-energy}

\clearpage
\section{Far Detector}

\fourfig{energy_fhc_fd_nonswap.png}{energy_fhc_fd_fluxswap.png}
        {energy_rhc_fd_nonswap.png}{energy_rhc_fd_fluxswap.png}
{FD reconstructed $\nu_e$ energy: non-swap (beam) left, flux-swap (appearance
signal) right; FHC top, RHC bottom.}{fig:fd-energy}

\fourfig{res_fhc_fd_nonswap.png}{res_fhc_fd_fluxswap.png}
        {res_rhc_fd_nonswap.png}{res_rhc_fd_fluxswap.png}
{FD fractional resolution $(E_{\rm reco}-E_{\rm true})/E_{\rm true}$, same layout
as Fig.~\ref{fig:fd-energy}: non-swap left, flux-swap right; FHC top, RHC bottom.
The flux-swap samples are visibly narrower, as expected for a pure signal $\nu_e$
sample against one containing NC $\pi^0$ background.}{fig:fd-resolution}

\twofig{ts_sp_fhc_fd_nonswap.png}{ts_sp_rhc_fd_nonswap.png}
{FD non-swap composition, FHC left and RHC right, again on the Spline axis only.
The beam sample is background-rich, with the same mix as the ND.}
{fig:fd-comp-nonswap}

\twofig{ts_sp_fhc_fd_fluxswap.png}{ts_sp_rhc_fd_fluxswap.png}
{FD flux-swap composition, FHC left and RHC right. The appearance-signal sample
is almost entirely $\nu_e$~CC by construction --- the flux-swap replaces the
$\nu_\mu$ flux with $\nu_e$, leaving almost nothing to be mistaken for signal.}
{fig:fd-comp-fluxswap}

\threefig{pidstack_fhc_fd_nonswap_CVNeLow.png}{pidstack_fhc_fd_nonswap_CVNeHigh.png}{pidstack_fhc_fd_nonswap_Peripheral.png}
{Composition per PID bin, FHC FD non-swap, with the purities in the legends:
$16.9\%$ (CVNe-Low), $74.2\%$ (CVNe-High) and $43.7\%$ (Peripheral) $\nu_e$~CC.
Purity rises with PID score in every sample; the peripheral sample falls between
the two core bins because it is uncontained, so it accepts events the core
rejects.}{fig:fd-pidstack}

\threefig{pidstack_fhc_fd_fluxswap_CVNeLow.png}{pidstack_fhc_fd_fluxswap_CVNeHigh.png}{pidstack_fhc_fd_fluxswap_Peripheral.png}
{Composition per PID bin, FHC FD flux-swap. Every bin is near-pure, but the trend
survives: the loose CVNe-Low bin still carries the most background, a couple of
percent NC, and the tight CVNe-High bin is the purest.}{fig:fd-pidstack-fluxswap}

\section{A Complete List of Relevant Docdbs with Citations Including URL}

\bibliographystyle{unsrtnat}
\bibliography{summary-ssahoo}

\end{document}
